\appendix
\section{Appendix}
\label{app:supplement}

\subsection{Data and code repository}
\label{app:hybrid-stats}

\noindent Mann--Whitney outputs under \path{sim/multi_seed/stats/}, when present, are optional regression checks after a CSV refresh. They are \textbf{not cited in this manuscript} (descriptive panels only in the main text).

\subsection{Learning-stratification sanity check (Fig.~\ref{fig:eval-learn})}
\label{app:learn}

\noindent\textbf{Scope.} Figure~\ref{fig:eval-learn} does \textbf{not} support causal or performance claims about bounded Q-learning. Section~\ref{sec:ql} is a \textbf{design specification only}.

The panel compares \texttt{learning\_on} strata exported from the \emph{same} \texttt{AER-MQoS} binary (\texttt{learn\_or\_load.csv}). It verifies that the logging path can emit both labels; it is \textbf{not} an independent firmware ablation (no separate \texttt{NO\_QL} build). Causal evaluation requires ablation~(A2) in Section~\ref{sec:validation-plan}.

\begin{figure}[ht]
  \centering
  \includegraphics[width=\linewidth]{Fig_11_Learning_Or_Load_Sensitivity}
  \caption{\CaptionFigEleven}
  \label{fig:eval-learn}

  \noindent Fig.~\ref{fig:eval-learn} compares PDR across learning modes, verifying that the Q-learning integration does not degrade basic forwarding performance.
\end{figure}
